\section{Day 36: Splitting Fields (Feb.\ 13, 2026)}
Our goal is to prove part (ii) of the theorem. Our subgoal is to show that, given $p \in F[x]$, there exists a smallest field containing all the roots of $p$.
\begin{claim}
    $a$ is a root of $p \in F[x]$ if and only if $(x - a) \mid p$.
\end{claim}
\begin{proof}
    $p = (x - a) q + r$, where $p(a) = r$ implies $r = 0$, so $(x - a) \mid p$.
\end{proof}
\begin{definition}
    $a$ is a root of $p$ of multiplicity $k$ if $(x - a)^k \mid p$, but $(x - a)^{k+1} \nmid p$.
\end{definition}
\begin{proposition}
    If $\deg p = n$, then $p$ may have at most $n$ roots counted with multiplicity.
\end{proposition}
\begin{proof}
    Given distinct roots $a_i$ with multiplicity $k_i$, $(x - a_i)^{k_i} \mid p$ implies that $\prod (x - a_i)^{k_i} \mid p$ implies the degree of the LHS is at most $\deg p$ as desired.
\end{proof}
\begin{theorem}
    If $p \in F[x]$ is a polynomial, then there exists an extension $E$ of $F$ such that $P$ has a root $a \in E$. Furthermore, $\coor{E:F} < \deg p$.
\end{theorem}
\begin{proof}
    Let $g$ be an irreducible polynomial in $F[x]$ such that $g \mid p$. Let $E = F[x]\,/\!\left<q\right>$ and $a = [x]$ an equivalence class, then $g(a) = g([x]) = [g(x)] = 0$, implying $a$ is a root of $g$, and hence of $p$.
\end{proof}
\begin{corollary}
    If $p \in F(x)$ is such that $\deg p = n$, then there is $E/F$ such that $p$ factors into linear terms in $E$.
    \[ p = c \prod_{c, a_i \in E} (x - a_i). \]
    Furthermore, $\coor{E:F} \leq n!$.
\end{corollary}
\begin{proof}
    Using the previous theorem, find $a_1 \in E_1 / F$ such that $p(a_1) = 0$, so $(x - a_1) \mid p$, and $p = (x - a_1) p_1$ with $\deg(p_1) = n-1$. Inductively, we find $E/E_1$ such that $p_1$ splits into a product of linear factors in $E$, and so $p$ splits in $E$. We can write
    \[ \coor{E:F} = \coor{E:E_1} \coor{E_1:F} \leq (n-1)! n = n!. \qedhere \]
\end{proof}
As an example, $p = x^3 - 2 \in Q[x]$. $E_1 = \QQ[x] \,/\! \left<x^3 - 2\right> = \QQ(2^{1/3})$. We have that $\coor{E_1:\QQ} = 3$ with roots $1, 2^{1/3}, 2^{2/3}$, but this does not fully split $x^3 - 2$. $p = (x - 2^{1/3}) p_1$, where 
\[ p_1 = \frac{x^3 - 2}{x - 2^{1/3}} = x^2 + 2^{1/3} x + 2^{2/3} \]
is irreducible and only has complex roots. Now, $E = E_1 \,/\! \left<x^2 + 2^{1/3} x + 2^{2/3}\right>$ is a field, in which $x^2 + 2^{1/3} x + 2^{2/3}$ has a factor, so it has two factors and thus splits, which we could compute with the quadratic formula if we wanted. Either way, we do have
\[ \coor{E:\QQ} = \coor{E:E_1} \coor{E_1:\QQ} = 2 \cdot 3 = 6. \] \vspace{-20pt}
\begin{definition}
    Let $p \in F[x]$. A splitting field $E$ for $p$ over $F$ is an extension $E/F$ such that $p$ factors into linear terms in $E$, and $p$ does not factor into linear terms in any proper subfield of $E$.
\end{definition}

\newpage
\begin{claim}
    If $p$ splits in $K$, $p = \prod (x - a_i)$ where $a_i \in K$, then $F(a_1, \dots, a_n)$ is a splitting field of $p$. It is, indeed, the smallest by definition of $F(\bullet)$.
\end{claim}
\begin{proof}
    Obviously, $p$ splits in $F(a_1, \dots, a_n)$.
\end{proof}
\begin{theorem}
    If $E$ and $E'$ are splitting fields of $f \in F[x]/F$, then we have
    \[ \begin{tikzcd} E \arrow[rr, "\sim"] \arrow[rr, "\phi"'] & & E' \\ & F \arrow[ul, hookrightarrow] \arrow[ur, hookrightarrow] & \end{tikzcd} \]
\end{theorem}